% GENERATED FILE. Edit canonical lessons, then run npm run generate:books.
% canonical-source-hash: fnv1a64:78e5d4cfb6f1e6cf
% canonical-lessons: TA-C349-cinti, TA-C349-uki, TA-C349-anumani, TA-C349-karutukol, TA-C349-palar

\chapter{Thinking It Through}
\label{ch:ta-thinking-it-through}
\hlchaptermodality{349}

\begin{chapteropening}
\textbf{By the end of this chapter:} \emph{I can say that I guess, assume or think something over.}

Complete the last lesson of chapter 349: I can say that I guess, assume or think something over.
\end{chapteropening}

\section[\texorpdfstring{\ta{சிந்தி} (cinti)}{சிந்தி (cinti)}]{\ta{சிந்தி} (cinti) — to think deeply; to reflect}

\label{lesson:TA-C349-cinti}



\begin{quote}
Before the new one: say the Tamil for a consensus, then the Tamil for some people.
\end{quote}

\subsection*{You'll want to know: \ta{சிந்தி}}
\textbf{\ta{சிந்தி}} — \emph{cinti} — \textquotedblleft{}to think deeply; to reflect\textquotedblright{}.

\begin{quote}
\ta{முடிவு} \ta{செய்வதற்கு} \ta{முன்} \ta{அவள்} \ta{நன்றாகச்} \ta{சிந்தித்தாள்}.
\end{quote}

\begin{quote}
\emph{muṭivu ceyvataṟku muṉ avaḷ naṉṟākac cintittāḷ.}
\end{quote}

\begin{quote}
She thought it over carefully before deciding.
\end{quote}

Past (I …): \ta{சிந்தித்தேன்}, \emph{cintittēṉ}.

What you think about takes -\ta{ஐப்} \ta{பற்றி}: \ta{அதைப்} \ta{பற்றிச்} \ta{சிந்தி}. The noun is \ta{சிந்தனை}, a thought.

\begin{cousinweb}
\textbf{\ta{சிந்தி}} (\emph{cinti}) is a verb built on Sanskrit \emph{cintā}, \textquotedblleft{}thought, care\textquotedblright{} (CDIAL 4816), with the loan-verb ending -\ta{இ}.
\end{cousinweb}

\begin{grammarlens}[title={What you've built}]
One more word for giving reasons and opinions.
\end{grammarlens}

\subsection*{Your turn}
\begin{itemize}
\raggedright
  \item \emph{Say it:} \emph{cinti}
  \item \emph{Say it:} \emph{cinti}, once more
  \item \emph{Say it:} \emph{muṭivu ceyvataṟku muṉ avaḷ naṉṟākac cintittāḷ.} — She thought it over carefully before deciding.
  \item \emph{Recall:} say the Tamil for a consensus, then the Tamil for some people, then say \emph{cinti} again
\end{itemize}

\begin{tcolorbox}[breakable,colback=teal!4,colframe=teal!35!black,title={Before you move on}]
What does \ta{சிந்தி} mean? (\textquotedblleft{}To think deeply; to reflect\textquotedblright{}.) Say it once more.
\end{tcolorbox}
\section[\texorpdfstring{\ta{ஊகி} (ūki)}{ஊகி (ūki)}]{\ta{ஊகி} (ūki) — to guess; to infer}

\label{lesson:TA-C349-uki}



\begin{quote}
Before the new one: say the Tamil for some people, then the Tamil for to think deeply.
\end{quote}

\subsection*{You'll want to know: \ta{ஊகி}}
\textbf{\ta{ஊகி}} — \emph{ūki} — \textquotedblleft{}to guess; to infer\textquotedblright{}.

\begin{quote}
\ta{அவன்} \ta{பதிலைச்} \ta{சரியாக} \ta{ஊகித்தான்}.
\end{quote}

\begin{quote}
\emph{avaṉ patilaic cariyāka ūkittāṉ.}
\end{quote}

\begin{quote}
He guessed the answer correctly.
\end{quote}

Past (I …): \ta{ஊகித்தேன்}, \emph{ūkittēṉ}.

What is guessed takes -\ta{ஐ}: \ta{பதிலை} \ta{ஊகி}, to guess the answer.

\begin{cousinweb}
\textbf{\ta{ஊகி}} (\emph{ūki}) is a verb built on Sanskrit \emph{ūha}, \textquotedblleft{}consideration\textquotedblright{}, Prakrit \textquotedblleft{}reasoning, guess\textquotedblright{} (CDIAL 2443), with the loan-verb ending -\ta{இ}.
\end{cousinweb}

\begin{grammarlens}[title={What you've built}]
One more word for giving reasons and opinions.
\end{grammarlens}

\subsection*{Your turn}
\begin{itemize}
\raggedright
  \item \emph{Say it:} \emph{ūki}
  \item \emph{Say it:} \emph{ūki}, once more
  \item \emph{Say it:} \emph{avaṉ patilaic cariyāka ūkittāṉ.} — He guessed the answer correctly.
  \item \emph{Recall:} say the Tamil for some people, then the Tamil for to think deeply, then say \emph{ūki} again
\end{itemize}

\begin{tcolorbox}[breakable,colback=teal!4,colframe=teal!35!black,title={Before you move on}]
What does \ta{ஊகி} mean? (\textquotedblleft{}To guess; to infer\textquotedblright{}.) Say it once more.
\end{tcolorbox}
\section[\texorpdfstring{\ta{அனுமானி} (aṉumāṉi)}{அனுமானி (aṉumāṉi)}]{\ta{அனுமானி} (aṉumāṉi) — to assume; to presume}

\label{lesson:TA-C349-anumani}



\begin{quote}
Before the new one: say the Tamil for to think deeply, then the Tamil for to guess.
\end{quote}

\subsection*{You'll want to know: \ta{அனுமானி}}
\textbf{\ta{அனுமானி}} — \emph{aṉumāṉi} — \textquotedblleft{}to assume; to presume\textquotedblright{}.

\begin{quote}
\ta{அவர்} \ta{வரமாட்டார்} \ta{என்று} \ta{நான்} \ta{அனுமானித்தேன்}.
\end{quote}

\begin{quote}
\emph{avar varamāṭṭār eṉṟu nāṉ aṉumāṉittēṉ.}
\end{quote}

\begin{quote}
I assumed that he would not come.
\end{quote}

Past (I …): \ta{அனுமானித்தேன்}, \emph{aṉumāṉittēṉ}.

The assumption is a clause with \ta{என்று}: … \ta{என்று} \ta{அனுமானி}.

\begin{cousinweb}
\textbf{\ta{அனுமானி}} (\emph{aṉumāṉi}) is a verb built on Sanskrit \emph{anumāna}, \textquotedblleft{}inference, thought\textquotedblright{} (CDIAL 328), the term Indian logic used for reasoning from a sign to what it shows.
\end{cousinweb}

\begin{grammarlens}[title={What you've built}]
One more word for giving reasons and opinions.
\end{grammarlens}

\subsection*{Your turn}
\begin{itemize}
\raggedright
  \item \emph{Say it:} \emph{aṉumāṉi}
  \item \emph{Say it:} \emph{aṉumāṉi}, once more
  \item \emph{Say it:} \emph{avar varamāṭṭār eṉṟu nāṉ aṉumāṉittēṉ.} — I assumed that he would not come.
  \item \emph{Recall:} say the Tamil for to think deeply, then the Tamil for to guess, then say \emph{aṉumāṉi} again
\end{itemize}

\begin{tcolorbox}[breakable,colback=teal!4,colframe=teal!35!black,title={Before you move on}]
What does \ta{அனுமானி} mean? (\textquotedblleft{}To assume; to presume\textquotedblright{}.) Say it once more.
\end{tcolorbox}
\section[\texorpdfstring{\ta{கருதுகோள்} (karutukōḷ)}{கருதுகோள் (karutukōḷ)}]{\ta{கருதுகோள்} (karutukōḷ) — a hypothesis}

\label{lesson:TA-C349-karutukol}



\begin{quote}
Before the new one: say the Tamil for to guess, then the Tamil for to assume.
\end{quote}

\subsection*{You'll want to know: \ta{கருதுகோள்}}
\textbf{\ta{கருதுகோள்}} — \emph{karutukōḷ} — \textquotedblleft{}a hypothesis\textquotedblright{}.

\begin{quote}
\ta{விஞ்ஞானி} \ta{ஒரு} \ta{புதிய} \ta{கருதுகோளை} \ta{முன்வைத்தார்}.
\end{quote}

\begin{quote}
\emph{viññāṉi oru putiya karutukōḷai muṉvaittār.}
\end{quote}

\begin{quote}
The scientist put forward a new hypothesis.
\end{quote}

\ta{கருதுகோள்} \ta{முன்வை}, to put forward a hypothesis.

\begin{cousinweb}
\textbf{\ta{கருதுகோள்}} (\emph{karutukōḷ}) is a modern Tamil coinage built from native \emph{karutu}, \textquotedblleft{}to intend, suppose, consider\textquotedblright{} (DEDR 1283), and \emph{kōḷ}, \textquotedblleft{}taking, opinion, tenet\textquotedblright{}, the noun of \ta{கொள்}: \textquotedblleft{}what is supposed\textquotedblright{}.
\end{cousinweb}

\begin{grammarlens}[title={What you've built}]
One more word for giving reasons and opinions.
\end{grammarlens}

\subsection*{Your turn}
\begin{itemize}
\raggedright
  \item \emph{Say it:} \emph{karutukōḷ}
  \item \emph{Say it:} \emph{karutukōḷ}, once more
  \item \emph{Say it:} \emph{viññāṉi oru putiya karutukōḷai muṉvaittār.} — The scientist put forward a new hypothesis.
  \item \emph{Recall:} say the Tamil for to guess, then the Tamil for to assume, then say \emph{karutukōḷ} again
\end{itemize}

\begin{tcolorbox}[breakable,colback=teal!4,colframe=teal!35!black,title={Before you move on}]
What does \ta{கருதுகோள்} mean? (\textquotedblleft{}A hypothesis\textquotedblright{}.) Say it once more.
\end{tcolorbox}
\section[\texorpdfstring{\ta{பலர்} (palar)}{பலர் (palar)}]{\ta{பலர்} (palar) — many people}

\label{lesson:TA-C349-palar}



\begin{quote}
Before the new one: say the Tamil for to assume, then the Tamil for a hypothesis.
\end{quote}

\subsection*{You'll want to know: \ta{பலர்}}
\textbf{\ta{பலர்}} — \emph{palar} — \textquotedblleft{}many people\textquotedblright{}.

\begin{quote}
\ta{பலர்} \ta{இந்தக்} \ta{கருத்தை} \ta{ஏற்றுக்கொண்டார்கள்}.
\end{quote}

\begin{quote}
\emph{palar intak karuttai ēṟṟukkoṇṭārkaḷ.}
\end{quote}

\begin{quote}
Many people accepted this view.
\end{quote}

\ta{பலருக்கு}, to many people: the dative of a word in -\ta{அர்} adds -\ta{உக்கு}.

\begin{cousinweb}
\textbf{\ta{பலர்}} (\emph{palar}) is native Dravidian: DEDR 3987 lists Tamil \emph{pala}, \textquotedblleft{}many, several\textquotedblright{}, and \emph{palar}, \textquotedblleft{}many persons\textquotedblright{}, with the personal plural -\ta{அர்}.
\end{cousinweb}

\begin{grammarlens}[title={What you've built}]
One more word for giving reasons and opinions.
\end{grammarlens}

\subsection*{Your turn}
\begin{itemize}
\raggedright
  \item \emph{Say it:} \emph{palar}
  \item \emph{Say it:} \emph{palar}, once more
  \item \emph{Say it:} \emph{palar intak karuttai ēṟṟukkoṇṭārkaḷ.} — Many people accepted this view.
  \item \emph{Recall:} say the Tamil for to assume, then the Tamil for a hypothesis, then say \emph{palar} again
\end{itemize}

\begin{tcolorbox}[breakable,colback=teal!4,colframe=teal!35!black,title={Before you move on}]
What does \ta{பலர்} mean? (\textquotedblleft{}Many people\textquotedblright{}.) Say it once more.
\end{tcolorbox}
